\documentclass[11pt]{article}
\usepackage[a4paper,margin=1in]{geometry}
\usepackage{fontspec}
\usepackage[boldfont=false]{xeCJK}
\setCJKmainfont{FandolSong-Regular.otf}
\usepackage{amsmath}
\usepackage{amssymb}
\usepackage{array}
\usepackage{booktabs}
\usepackage{graphicx}
\usepackage{adjustbox}
\usepackage{enumitem}
\setlist[itemize]{topsep=2pt,partopsep=0pt,parsep=0pt,itemsep=2pt,leftmargin=1.4em}
\setlist[enumerate]{topsep=2pt,partopsep=0pt,parsep=0pt,itemsep=2pt,leftmargin=1.6em}
\usepackage{parskip}
\usepackage{fvextra}
\fvset{breaklines=true,breakanywhere=true,fontsize=\small}
\setlength{\parindent}{0pt}

\begin{document}

\begin{center}
  {\LARGE\bfseries Software Development}\\[0.35em]
  {\large A Level Computer Science}
\end{center}
\vspace{0.6em}

\section*{Program development life cycle}

A \textbf{development life cycle} 开发生命周期 is the set of stages from idea to finished, maintained software. It exists to \textbf{plan, manage and control} a project — to build the right product, on time, with good quality.

\subsection*{Why a life cycle is needed}

It manages complexity (break a big program into phases), coordinates teams, tracks progress with milestones, builds in testing, records design decisions for later, and manages risk.

\subsection*{Why there are different ones}

No single life cycle fits every project. The choice depends on the size and complexity, how clear the \textbf{requirements} 需求 are at the start, how much change is expected, the risk level, the team, and the deadline.

\subsection*{Common models}

\begin{itemize}
\item \textbf{Waterfall} 瀑布模型 — a linear sequence (Analysis → Design → Coding → Testing → Maintenance), each stage finished before the next. Clear and well-documented; good for \textbf{stable requirements}, but poor at coping with mid-project change, and the customer sees nothing working until the end.

\item \textbf{Iterative model} 迭代模型 — repeated passes, each producing a partial version that is reviewed and refined. Catches problems earlier; good when requirements are \textbf{discovered over time}, but harder to estimate.

\item \textbf{Rapid Application Development} 快速应用开发 (RAD) — heavy use of a \textbf{prototype} 原型 and user feedback. Very fast first delivery; good for \textbf{changing requirements}, but depends on user availability and suits smaller systems.

\item \textbf{Agile} 敏捷 — short iterations ("sprints"), constant collaboration and testing. Flexible and adaptive, but needs a committed customer and a skilled team.

\end{itemize}

\par\medskip
\noindent\begin{minipage}{\linewidth}
\centering
\adjustbox{max width=\linewidth, max totalheight=0.42\textheight}{%
  \includegraphics{/Users/joschka/Desktop/Programming/eXercise/output/eXam/handouts/a_level_computer_science/assets/12-p301-waterfall-model.pdf}}\\[0.3em]
\emph{The waterfall model: each stage is finished before the next begins}
\end{minipage}\par\medskip

\par\medskip
\noindent\begin{minipage}{\linewidth}
\centering
\adjustbox{max width=\linewidth, max totalheight=0.42\textheight}{%
  \includegraphics{/Users/joschka/Desktop/Programming/eXercise/output/eXam/handouts/a_level_computer_science/assets/12-p302-iterative-model.pdf}}\\[0.3em]
\emph{The iterative model: repeated passes refine the program}
\end{minipage}\par\medskip

\par\medskip
\noindent\begin{minipage}{\linewidth}
\centering
\adjustbox{max width=\linewidth, max totalheight=0.42\textheight}{%
  \includegraphics{/Users/joschka/Desktop/Programming/eXercise/output/eXam/handouts/a_level_computer_science/assets/12-p302-rad-model.pdf}}\\[0.3em]
\emph{Rapid application development: teams work on parts in parallel}
\end{minipage}\par\medskip

\subsection*{The standard stages}

\begin{itemize}
\item \textbf{analysis} — find \textbf{what} the program must do; gather and document requirements.

\item \textbf{design} — decide \textbf{how}: data structures, algorithms, modules, interface, file layouts.

\item \textbf{coding (implementation} 实现**)** — write the source code following the design.

\item \textbf{testing} — run against test data and fix bugs.

\item \textbf{maintenance} 维护 — after release, keep it working and useful.

\end{itemize}

\section*{Program design tools}

\subsection*{Structure chart}

A \textbf{structure chart} 结构图 shows the \textbf{hierarchical decomposition} 分解 of a program into modules (\textbf{subroutines} 子程序) and the \textbf{parameters} 参数 passed between them. Each module is a rectangle; lines link caller (above) to callee (below); small arrows show data going down and results coming back up.

\begin{Verbatim}
                CalculatePay
            /        |         \
       GetEmployee  CalculateBonus  CalculateTax
       Returns:     Takes: sales    Takes: gross
       employeeID   Returns: bonus  Returns: tax
\end{Verbatim}

It is a design-stage tool, and you can read the procedure signatures off it.

\par\medskip
\noindent\begin{minipage}{\linewidth}
\centering
\adjustbox{max width=\linewidth, max totalheight=0.42\textheight}{%
  \includegraphics{/Users/joschka/Desktop/Programming/eXercise/output/eXam/handouts/a_level_computer_science/assets/12-p304-structure-chart.pdf}}\\[0.3em]
\emph{A structure chart: modules with the parameters passed between them}
\end{minipage}\par\medskip

\subsection*{State-transition diagram}

A \textbf{state-transition diagram} 状态转换图 shows the \textbf{states} 状态 a system can be in and the \textbf{events} that move it between them — good for vending machines, traffic lights, user interfaces. Each state is a circle; each transition is an arrow labelled with the event.

\begin{Verbatim}
   coin inserted               item selected
[Idle] ──────────────→ [Awaiting selection] ──────────→ [Dispensing]
\end{Verbatim}

It makes missing transitions easy to spot ("what if a second coin is inserted while awaiting selection?").

\par\medskip
\noindent\begin{minipage}{\linewidth}
\centering
\adjustbox{max width=\linewidth, max totalheight=0.42\textheight}{%
  \includegraphics{/Users/joschka/Desktop/Programming/eXercise/output/eXam/handouts/a_level_computer_science/assets/12-p309-state-transition.pdf}}\\[0.3em]
\emph{A state-transition diagram for a door lock with code 259}
\end{minipage}\par\medskip

\section*{Errors}

\begin{itemize}
\item \textbf{syntax error} 语法错误 — breaks the language's grammar (missing bracket, misspelled keyword). Caught at translation time; the program won't run until fixed.

\item \textbf{run-time error} 运行时错误 — happens while running (divide by zero, file not found, array index out of range). The program crashes or raises an exception; fix by adding checks.

\item \textbf{logic error} 逻辑错误 — the program runs but gives \textbf{wrong results} (using \texttt{+} for \texttt{-}, an off-by-one loop, conditions in the wrong order). The hardest to find; the only sign is wrong output, so use careful testing and tracing.

\end{itemize}

\section*{Testing methods}

\begin{itemize}
\item \textbf{dry run} 手工跟踪 — trace the code on paper, writing each variable's value in a table.

\item \textbf{walkthrough} 走查 — a team review of the code.

\item \textbf{white-box testing} 白盒测试 — designed from the code's internal structure, covering every statement, branch and loop.

\item \textbf{black-box testing} 黑盒测试 — designed from the specification only: feed inputs, check outputs.

\item \textbf{integration testing} 集成测试 — combine modules and test the interfaces between them.

\item \textbf{alpha testing} α测试 — by the developers/in-house before release; \textbf{beta testing} β测试 — by a limited group of real users in their own environment.

\item \textbf{acceptance testing} 验收测试 — by the customer, to decide if the product is fit for purpose.

\item \textbf{stub} 桩 — a placeholder for a module that does not exist yet, so the structure can be tested top-down.

\end{itemize}

\section*{Test strategy and test plan}

A \textbf{test strategy} 测试策略 is the \textbf{high-level approach} — which kinds of testing, who does them, when, and the criteria to move on. A \textbf{test plan} 测试计划 is the \textbf{detailed list of tests} — each with input data, expected output, and a column for the actual output.

\subsection*{Choosing test data}

For each field or condition, include three kinds:

\begin{itemize}
\item \textbf{normal data} 正常数据 — typical values inside the valid range (for marks 0–100: \texttt{50}, \texttt{75}).

\item \textbf{abnormal data} 异常数据 — values that should be rejected (\texttt{-10}, \texttt{200}, \texttt{"abc"}).

\item \textbf{boundary data} 边界数据 — values at the edges, where off-by-one errors hide (\texttt{0}, \texttt{100}, and just outside \texttt{-1}, \texttt{101}).

\end{itemize}

\section*{Maintenance}

Most of a program's lifetime cost is in maintenance. Three kinds:

\begin{itemize}
\item \textbf{perfective maintenance} 完善性维护 — improving performance or features even though it works (a faster query, a new option).

\item \textbf{adaptive maintenance} 适应性维护 — keeping it working in a changing environment (a new OS, a new API, a legal change).

\item \textbf{corrective maintenance} 纠正性维护 — fixing bugs found in use.

\end{itemize}

A program may need all three throughout its life.

\section*{Amending an existing program}

When asked to add a feature or fix a bug:

\begin{enumerate}
\item \textbf{read the existing code} until you understand the algorithm and data flow.

\item \textbf{find where the change goes} — which subroutine, which lines.

\item \textbf{make the change as small as possible} — don't rewrite working code.

\item \textbf{update related parts} — every caller of a changed parameter list, every routine using a changed data structure.

\item \textbf{test} the new behaviour \textbf{and} the old (\textbf{regression testing} 回归测试 — check you broke nothing).

\item \textbf{document} the change.

\end{enumerate}

Clear comments, meaningful names, decomposed subroutines and a structure chart make a program much easier to amend — which is why the design tools matter even after the first release.


\end{document}
